\section*{Sets and Indices}

\begin{itemize}
    \item $T$: set of toy types, indexed by $t$. In this instance,
    \[
    T=\{\text{High-End},\ \text{Mid-Range},\ \text{Low-End}\}.
    \]
\end{itemize}

\section*{Parameters}

For each toy type $t \in T$:
\begin{itemize}
    \item $a_t$: manufacturing labor hours required per unit.
    \item $b_t$: inspection hours required per unit.
    \item $p_t$: profit per unit.
    \item $d_t$: maximum demand.
\end{itemize}

Global resource limits:
\begin{itemize}
    \item $L$: total available labor hours.
    \item $I$: total available inspection hours.
\end{itemize}

Using the CSV data for this instance:
\[
L=1000, \qquad I=500.
\]

Parameter values by toy type:
\[
\begin{array}{c|cccc}
t & a_t & b_t & p_t & d_t \\ \hline
\text{High-End} & 17 & 8 & 300 & 50 \\
\text{Mid-Range} & 10 & 4 & 200 & 80 \\
\text{Low-End} & 2 & 2 & 100 & 150
\end{array}
\]

\section*{Decision Variables}

For each toy type $t \in T$:
\begin{itemize}
    \item $x_t$ (code: \texttt{x[toy['type']]}): production quantity of toy type $t$.
\end{itemize}

The code defines these variables as continuous and nonnegative:
\[
x_t \ge 0 \qquad \forall t \in T.
\]

\section*{Objective}

Maximize total profit:
\[
\max \sum_{t \in T} p_t x_t.
\]

For this instance:
\[
\max\ 300x_{\text{High-End}} + 200x_{\text{Mid-Range}} + 100x_{\text{Low-End}}.
\]

\section*{Constraints}

Labor-hour capacity:
\[
\sum_{t \in T} a_t x_t \le L.
\]

Inspection-hour capacity:
\[
\sum_{t \in T} b_t x_t \le I.
\]

Demand upper bounds:
\[
x_t \le d_t \qquad \forall t \in T.
\]

Expanded for this instance:
\[
17x_{\text{High-End}} + 10x_{\text{Mid-Range}} + 2x_{\text{Low-End}} \le 1000,
\]
\[
8x_{\text{High-End}} + 4x_{\text{Mid-Range}} + 2x_{\text{Low-End}} \le 500,
\]
\[
x_{\text{High-End}} \le 50,\qquad
x_{\text{Mid-Range}} \le 80,\qquad
x_{\text{Low-End}} \le 150.
\]

\section*{Notes on Implementation}

\begin{itemize}
    \item The implemented model is a linear program with continuous production variables; no integrality restrictions are imposed.
    \item The code constructs one variable per toy type using the toy names read from \texttt{table\_1.csv}.
    \item Demand limits are read from \texttt{general\_parameters.csv} and matched by toy type.
    \item The model is solved using \texttt{GUROBI\_CMD(msg=0)} through a PuLP-compatible interface.
\end{itemize}
